
% \documentclass[handout]{beamer}[10]
\documentclass{beamer}[10] 
\usepackage{pgf}
\usepackage[english]{babel}
\usepackage[utf8]{inputenc}
\usepackage{beamerthemesplit}
\usepackage{graphics,epsfig, subfigure}
\usepackage{url}
\usepackage{srcltx}
\usepackage{hyperref}
\usepackage{bookmark}
\usepackage{nicefrac}
\usepackage{multicol}
\usepackage{vwcol}
\usepackage{microtype}
\usepackage[all]{xy}
\usepackage{amssymb}
\usepackage{tikz-cd}
% \usetikzlibrary{cd}


\definecolor{kugreen}{RGB}{50,93,61}
\definecolor{kugreenlys}{RGB}{132,158,139}
\definecolor{kugreenlyslys}{RGB}{173,190,177}
\definecolor{kugreenlyslyslys}{RGB}{214,223,216}
\setbeamercovered{invisible}
\mode<presentation>
\usetheme[numbers,totalnumber,compress,sidebarshades]{PaloAlto}
\setbeamertemplate{footline}[frame number]

  \usecolortheme[named=kugreen]{structure}
  \useinnertheme{circles}
  \usefonttheme[onlymath]{serif}
  \setbeamercovered{invisible}
  \setbeamertemplate{blocks}[rounded][shadow=true]

% \logo{\colorbox{kugreen}{\strut \hspace{0.37cm}\includegraphics[width=0.7cm]{HUJI_LOGO}\hspace{0.25cm}}}
\beamertemplatenavigationsymbolsempty
%\useoutertheme{infolines} 
\title{Equations in Hyperbolic(-esque) Groups}
\author{Barak Ohana}
\institute{The Hebrew University of Jerusalem}
\date{9 Sep 2024}

\newcommand{\sol}{\text{V}_{\Gamma}}
\renewcommand{\hom}{\text{Hom}}

%mathcal

\newcommand{\cala}{\mathcal{A}}
\newcommand{\calb}{\mathcal{B}}
\newcommand{\calc}{\mathcal{C}}
\newcommand{\cald}{\mathcal{D}}
\newcommand{\cale}{\mathcal{E}}
\newcommand{\calf}{\mathcal{F}}
\newcommand{\calg}{\mathcal{G}}
\newcommand{\calh}{\mathcal{H}}
\newcommand{\cali}{\mathcal{I}}
\newcommand{\calj}{\mathcal{J}}
\newcommand{\calk}{\mathcal{K}}
\newcommand{\call}{\mathcal{L}}
\newcommand{\calm}{\mathcal{M}}
\newcommand{\caln}{\mathcal{N}}
\newcommand{\calo}{\mathcal{O}}
\newcommand{\calp}{\mathcal{P}}
\newcommand{\calq}{\mathcal{Q}}
\newcommand{\calr}{\mathcal{R}}
\newcommand{\cals}{\mathcal{S}}
\newcommand{\calt}{\mathcal{T}}
\newcommand{\calu}{\mathcal{U}}
\newcommand{\calv}{\mathcal{V}}
\newcommand{\calw}{\mathcal{W}}
\newcommand{\calx}{\mathcal{X}}
\newcommand{\caly}{\mathcal{Y}}
\newcommand{\calz}{\mathcal{Z}}

%mathbb
\usepackage{amsfonts}
\renewcommand{\aa}{\mathbb{A}}
\newcommand{\bb}{\mathbb{B}}
\newcommand{\cc}{\mathbb{C}}
\newcommand{\dd}{\mathbb{D}}
\newcommand{\ee}{\mathbb{E}}
\newcommand{\ff}{\mathbb{F}}
\renewcommand{\gg}{\mathbb{G}}
\newcommand{\hh}{\mathbb{H}}
\newcommand{\ii}{\mathbb{I}}
\newcommand{\jj}{\mathbb{J}}
\newcommand{\kk}{\mathbb{K}}
\renewcommand{\ll}{\mathbb{L}}
\newcommand{\mm}{\mathbb{M}}
\newcommand{\nn}{\mathbb{N}}
\newcommand{\oo}{\mathbb{O}}
\newcommand{\pp}{\mathbb{P}}
\newcommand{\qq}{\mathbb{Q}}
\newcommand{\rr}{\mathbb{R}}
\renewcommand{\ss}{\mathbb{S}}
\renewcommand{\tt}{\mathbb{T}}
\newcommand{\uu}{\mathbb{U}}
\newcommand{\vv}{\mathbb{V}}
\newcommand{\ww}{\mathbb{W}}
\newcommand{\xx}{\mathbb{X}}
\newcommand{\yy}{\mathbb{Y}}
\newcommand{\zz}{\mathbb{Z}}

%parentheses
\newcommand*{\rb}[1]{\left( #1 \right)}
\renewcommand*{\sb}[1]{\left[ #1 \right]}
\newcommand*{\cb}[1]{\left\{ #1 \right\}}
\newcommand*{\ab}[1]{\left< #1 \right>}
\newcommand*{\nb}[1]{\lVert #1 \rVert}
\newcommand{\size}[1]{\left| #1 \right|}

\newcommand*{\by}{\pmb{Y}}
\DeclareMathOperator{\diam}{diam}
\DeclareMathOperator{\mcg}{MCG}
\newcommand{\sker}{\underset{\rightarrow}{\ker}}
\newcommand{\onto}{\twoheadrightarrow}
\DeclareMathOperator{\cay}{Cay}
\newcommand{\acts}{\curvearrowright}

\graphicspath{ {images/} }



\begin{document}

%frame 1
\begin{frame}\titlepage \vspace{-0.5cm}\end{frame}


\section{Equations in Groups}

%frame 2
\begin{frame}
    \frametitle{Equations in Groups}
    \pause
    \begin{itemize}
        \item $\Gamma$ - A group\pause
        \item $\Sigma\left(x_{1},\ldots,x_{d}\right)$ - A (possibly infinite) system of equation on d variables. \pause
    \end{itemize}

   
    We define the \textbf{set of solution of $\Sigma$ in $\Gamma$} to be the set
    \begin{equation*}
        \sol\left(\Sigma\right)=\left\{ \left(g_{1},\ldots,g_{d}\right)\in\Gamma^{d}\mid\Sigma\left(g_{1},\ldots,g_{d}\right)=_{\Gamma}1\right\} 
    \end{equation*}

            \pause

    There is an equivalence
    \begin{equation*}
        \sol\rb{\Sigma} \leftrightarrow \hom\rb{\nicefrac{\ff_d}{\ab{\ab{\Sigma}}}, \Gamma}
    \end{equation*}

    \pause

    So instead of studying solutions in $\Gamma$ we will study homomorphisms to $\Gamma$.

\end{frame}

%frame 3
\begin{frame}

    \frametitle{Equations in Groups}
    \pause
    \begin{definition}
        A group $\Gamma$ is called equationally noetherian if  every set of equation is equivalent to a finite subset of it. \pause
        
        i.e. if for any set of equation $\Sigma$ there exists a finite subset $\Sigma_0$ such that $\sol\rb{\Sigma} = \sol\rb{\Sigma_0}$.
    \end{definition}
    \pause
    \begin{example}\pause
        \begin{multicols}{2}
        \begin{itemize}
            \item All finite groups \pause
            \item All linear groups\pause
                \begin{itemize}
                    \item Thus all free groups\pause
                \end{itemize}
            \item Hyperbolic groups\\\pause
            \item[] 
            \item Free product of equationally noetherian groups\pause
            \item Groups hyperbolic relative to equationally noetherian
        \end{itemize}
        \end{multicols}
    \end{example}

\end{frame}

\section{Limit groups}

%frame 4
\begin{frame}
    \frametitle{Limit groups}
    \pause 
    \begin{definition}\pause
        Let $\Gamma$ be f.g. and $\rb{\varphi_n} \in \hom\rb{G,\Gamma}$ for some f.g. group $G$. \pause
        We define the \textbf{stable kernel} of $\rb{\varphi_n}$ to be
        \begin{equation*}
            \sker{\varphi_n} = \cb{g\in G \mid \varphi_n \rb{g} \text{  is eventually } 1}
        \end{equation*}\pause
        We define the corresponding \textbf{limit group} by $L=\nicefrac{G}{\sker{\varphi_n}}$
    \end{definition}
    \pause
    \begin{lemma}\pause
        Let $\Gamma$ be a group. The following are equivalent:\pause

        \begin{itemize}
            \item $\Gamma$ is equationally noetherian
            \pause \item For any f.g. $G$ and any $\rb{\varphi_n} \in \hom\rb{G,\Gamma}$, \pause $\varphi_n$ eventually factors through the
            quotient $G\onto L$
        \end{itemize}

    \end{lemma}
\end{frame}

%frame 5
\begin{frame}
    \frametitle{Limit groups}
    \begin{overprint}
        \onslide<2|handout:0>
            \begin{equation*}
                \xymatrix{
                    G\ar@/^12pt/[ddr]^{\varphi_n} \ar@{->>}[d] &  \\
                    L &  \\
                    & \Gamma
                }
            \end{equation*}
        \onslide<3->
            \begin{equation*}
                \xymatrix{
                    G\ar@/^12pt/[ddr]^{\varphi_n} \ar@{->>}[d] &  \\
                    L\ar@{-->}[dr]_{\hat{\varphi}_n} &  \\
                    & \Gamma
                }
            \end{equation*}
    \end{overprint}\pause\pause\pause 
    \begin{block}{Reminder}
        A single morphism $\varphi_n \in \hom\rb{G,\Gamma}$ factors through $G\onto L$ if and only if $\sker\varphi_n \subseteq \ker \varphi_n $. \pause
        
        Thus we want to show that we have that eventually we have $\sker\varphi_n \subseteq \ker \varphi_n $.
    \end{block}
\end{frame}

\section{Limit groups over Hyp. Groups}
%frame 6
\begin{frame} \frametitle{Limit groups over Hyperbolic Groups}

    \begin{itemize}
        \item Let $\Gamma$ be a f.g. hypebolic group (meaning its cayley graph is hyperbolic).
        \item \pause Let $\rb{\varphi_n} \in \hom\rb{G,\Gamma}$. 
        \item \pause Denote 
            \begin{equation*}
                \nb{\varphi_n} = \max_{s\in S} d_\Gamma \rb{1,\varphi_n\rb{ s }}
            \end{equation*}
            i.e. the maximal length of the images of the generators of $G$ under $\varphi_n$.\pause
    \end{itemize}
    
    
    \begin{fact}
        If $\nb{\varphi_n}$ is a bounded sequence then $\varphi_n$ has finitely many options (as morphisms).
        
        Thus if $\nb{\varphi_n}$ is a bounded sequence then $\varphi_n$ eventually factors through $G\onto L$.
    \end{fact}

\end{frame}

%frame 7
\begin{frame}
    \frametitle{Limit groups over Hyperbolic Groups}\pause
    \begin{block}{Proposition (Sela: torsion free hyp. groups, Weidmann-Reinfeldt: all hyp. groups)}
        $L$ is finitely presented relative to finitely many finitely generated subgroups $P_1, \ldots, P_d$ such that $\nb{\varphi_{n}\big\vert_{\tilde{P}_i}}$ is bounded for every $i$.
    \end{block}
    \pause
    {\Large What we have here?}
    \pause

    Finitely presented relative to $P_1, \ldots, P_d$ means that
    \begin{equation*}
        L = \ab{ P_1, \ldots, P_d, S \mid R} 
    \end{equation*}
    \pause where $S$ is a finite set of (extra) generators\pause, and $R$ is a finite set of (extra) relations.

\end{frame}

%frame 8
\begin{frame}
    \frametitle{Limit groups over Hyperbolic Groups}
    \vspace*{-3pt}
    \pause
    \begin{block}{Theorem (Sela, Weidmann-Reinfeldt)}
        If $\Gamma$ is a finitely generated hyperbolic group then $\Gamma$ is equationally noetherian.
    \end{block}
    \pause
    \vspace*{-3pt}
    \begin{proof}
        Let $\rb{\varphi_n} \in \hom\rb{G,\Gamma}$.\pause We want to show that $\varphi_n$ eventually factors through $G \onto L$.\pause

        Let $P_1,\ldots , P_d$ be the subgroups from before\pause, i.e $L$ is finitely presented relative to $P_1,\ldots , P_d$ and $\nb{\varphi_{n}\big\vert_{\tilde{P}_i}}$ are bounded sequences.\pause

        For each $i$, the sequence $\nb{\varphi_{n}\big\vert_{\tilde{P}_i}}$ is bounded thus $\varphi_{n}\big\vert_{\tilde{P}_i}$ has finitely many options as homoporphism.
        \pause This means that $\varphi_{n}\big\vert_{\tilde{P}_i} : \tilde{P}_i \to \Gamma$ factors through $\tilde{P}_i\onto P_i$. \pause Since $L$ is f.p. relative to the $P_i$'s then also $\varphi_n$ factors through $G\onto L$.

    \end{proof}
    % \vspace*{4ex}
    % But what happens with groups that are not hyperbolic?
\end{frame}

\section*{Hyperbolic-esque Groups}

%frame 9
\begin{frame}
    \frametitle{Hyperbolic-esque Groups}
    What can we say about non-hyperbolic groups?
    \pause

    \begin{itemize}
        \item Replace the assumption that $\cay\rb{ \Gamma }$ is hyperbolic by assuming that $\Gamma$ acts on hyperbolic space $X$.
        \item \pause Under some assumption on the action, the same Proposition as before apply.
    \end{itemize}
    \pause
    \begin{block}{Theorem (Groves \& Hull, 2017)}
        \pause Let $\Gamma$ be a group acting acylindrically on a hyperbolic space $X$. \pause  
        
        Let $L$ be a limit group defined by $\rb{\varphi_n} \in \hom\rb{G,\Gamma}$. \\ \pause
        
        $L$ is finitely presented relative to finitely many finitely generated subgroups $P_1, \ldots, P_d$ such that $\nb{\varphi_{n}\big\vert_{\tilde{P}_i}}$ is bounded for every $i$.
    \end{block}
    \pause


    Since $X$ is not the Cayley graph we dont know that if $\nb{\varphi_n}$ is bounded then $\varphi_n$ eventually factors through $G\onto L$.
\end{frame}

%frame 10
\begin{frame}
    \frametitle{Group acting on a Product of hyperbolic Spaces}
    {\Large Possible approach: Instead of having one action with assumptions, have two actions}\pause  
    \begin{block}{Proposition (O', 2023)}
        Let $\Gamma$ be a group acting acylindrically on two hyperbolic metric spaces, $X_1$ and $X_2$.\pause
            
        Assume that $\Gamma \acts X_1 \times X_2$ is like the actions of $\Gamma$ on $\cay\rb{\Gamma}$.\pause

        Under thouse assumption, $\Gamma$ is equationally noetherian
    \end{block}\pause

    \begin{block}{Proof.}\pause
        For a moment, look just on the action $\Gamma \acts X_1$.\pause 

        Let $P_1, \ldots, P_d$ be the groups from the Proposition before\pause, i.e. $L$ is f.p. relative to $P_1, \ldots, P_d$ and $\nb{\varphi_{n}\big\vert_{\tilde{P}_i}}_{X_1}$ (mesured in $X_1$!) is bounded.\pause

         We will explain why $\varphi_{n}\big\vert_{\tilde{P}_i}$ factors through $\tilde{P}_i \onto P_i$. 
    \end{block}
    
    
\end{frame}

%frame 11
\begin{frame}
    \frametitle{Group acting on a Product of hyperbolic Spaces}
    \begin{block}{Proof.}
        Restrict our attention to a specific $\tilde{P}_i \leq G$. \pause
        Look at $\varphi_{n}\big\vert_{\tilde{P}_i}: \tilde{P}_i \to \Gamma$ but now we focus at the action on $X_2$. \pause

        The Proposition before (applied to the action on $X_2$) tells us that $P_i$ is finitely presented relative to $Q_1, \ldots, Q_l$ where
        \begin{equation*}
            \nb{\varphi_{n}\big\vert_{\tilde{Q}_j}}_{X_2}
        \end{equation*}
        is bounded. Now mesured relative to $X_2$! \pause

        Since $\tilde{Q}_l$ is a subgroup of $P_i$, then both sequences
        \begin{equation*}
            \nb{\varphi_{n}\big\vert_{\tilde{Q}_j}}_{X_1};\hspace{1ex}\nb{\varphi_{n}\big\vert_{\tilde{Q}_j}}_{X_2}
        \end{equation*}
        are bounded. 
    \end{block}
\end{frame}

%frame 12
\begin{frame}
    \frametitle{Group acting on a Product of hyperbolic Spaces}
    \begin{proof}
    
        \pause This means that 
        \begin{equation*}
            \nb{\varphi_{n}\big\vert_{\tilde{Q}_j}}_{X_1 \times X_2} = \max_{s\in S} d_{X_1 \times X_2} \rb{1, \varphi_n \rb{s}}
        \end{equation*} 
        is bounded! \pause

        We now remember that $X_1 \times X_2$ is the cayley graph of $\Gamma$, which means that $\varphi_{n}\big\vert_{\tilde{Q}_j}$ has finitely many options as a homorphism.\pause

        This means $\varphi_{n}\big\vert_{\tilde{Q}_j}$ eventually factors through $\tilde{Q}_j \onto Q_j$. \pause
        
        Since this is true for all $Q_1, \ldots, Q_l$, then, as before, $\varphi_{n}\big\vert_{\tilde{P}_i}$ factors through $\tilde{P}_i \onto P_i$.\pause

        Wrapping this up, and doing this for all $P_1, \ldots, P_d$ we get that $\varphi_n$ eventually factor through $G \onto L$

    \end{proof}
\end{frame}

%frame 13
\begin{frame}
    \frametitle{Group acting on a Product of hyperbolic Spaces}

    The Proposition works for product of any length, i.e. we have
    \pause
    \begin{block}{Theorem (O')}
        Every f.g. group acting  geometrically and strictly acylindrically (i.e as before) on a product of hyperbolic spaces is equationally noetherian.
    \end{block}\pause

    This method generalize to groups acting on colorable Hierarchically Hyperbolic spaces (whatever this means)
    \pause
    \begin{block}{Theorem (O')}
        Every f.g. group acting geometrically and strictly acylindrically on a colorable HHG is equationally noetherian.
    \end{block}
    \pause
    \begin{center}
        {\huge The end!}
    \end{center}


\end{frame}



\end{document}